\chapter*{ABSTRACT}
\addcontentsline{toc}{chapter}{ABSTRACT}

Air pollution, particularly concerning fine particulate matter (PM$_{2.5}$), has emerged as a critical environmental concern in Vietnam with direct implications for public health. Accurate short-term forecasting of PM$_{2.5}$ concentrations (1–24 hours) plays a pivotal role in early warning systems and decision-making support. This study develops and benchmarks \textbf{six advanced models} for multi-horizon PM$_{2.5}$ prediction (T+1, T+3, T+6, T+12, T+24 hours), namely: XGBoost, XLinear, ESTGCN, ST-XLinear, Ensemble, and iTransformer. The dataset was acquired from \textbf{12 sensor monitoring stations} (6 in Northern and 6 in Southern Vietnam), deeply integrating air quality metrics (PM$_{2.5}$, PM$_{10}$, CO, NO$_2$, SO$_2$, O$_3$) with meteorological sequences (temperature, relative humidity, wind dynamics, precipitation, and soil temperature). A comprehensive, 12-step preprocessing pipeline was devised, encompassing the detection of Optical Particle Counter (OPC) failures, mitigation of sensor freezing anomalies, outlier handling, and the synthesis of six novel engineered features grounded in atmospheric chemistry (e.g., oxidation potential, humid sulfate risk, and thermal stability). Notably, to effectively model the spatiotemporal propagation of pollution among stations, this research proposes a \textbf{Dynamic Graph Builder}—an adjacency matrix strictly modulated by real-time wind directions ($\alpha = 0.4858$). Experimental results demonstrate that \textbf{XGBoost achieves the highest predictive efficiency} consistently across all forecast horizons, yielding an $R^2 = 90.47\%$ (T+1) and $70.95\%$ (T+24), significantly outperforming the comparative deep learning architectures. The mean RMSE ranges precisely from 5.16 $\mu g/m^3$ (T+1) to 11.52 $\mu g/m^3$ (T+24). Furthermore, the Southern station cluster exhibited superior accuracy metrics compared to the Northern cluster ($R^2 = 98.81\%$ vs. $82.12\%$ at T+1), reflecting a lower intrinsic volatility of PM$_{2.5}$ particles attributable to the region-specific climatic wash-out effect. Ultimately, the entire algorithmic framework was successfully deployed into an interactive Streamlit web application, facilitating real-time monitoring and online forecasting of PM$_{2.5}$ across all 12 stations, providing a highly scalable and functional public health early warning system.